% II. Проектиране на архитектура на системата и базата данни
% Ориентировъчен обем: 10–11 стр.
\section{Проектиране на архитектура на системата и базата данни}
\label{sec:architecture}

\subsection{Критерии за избор на архитектура}
\label{subsec:criteria}

Изискванията от Раздел~\ref{sec:domain} не определят еднозначно нито едно от
основните архитектурни решения. За да бъде изборът обоснован, а не въпрос на вкус,
четирите критерия за оценка са записани преди разглеждането на вариантите:

\begin{itemize}[topsep=2pt,itemsep=2pt,leftmargin=*]
\item \textbf{К1, проверимост.} Може ли твърдението „системата дава верните числа“ да
      бъде проверено механично, върху целия входен обхват, без човек да сравнява
      клетки. Този критерий е поставен първи, защото NFR1 е единственото изискване,
      чието нарушаване прави всичко останало безполезно.
\item \textbf{К2, работоспособност при частичен отказ.} Какво остава използваемо,
      когато една от частите на системата не е достъпна (NFR2).
\item \textbf{К3, цена на доставянето.} Какво е необходимо, за да бъде публикувана нова
      версия: инструментална верига, стъпка на компилация, среда за изпълнение (NFR7).
\item \textbf{К4, повторна употреба между клиентите.} Каква част от изходния текст се
      използва без промяна и от уеб, и от мобилния клиент (NFR3).
\end{itemize}

\subsection{Разгледани архитектурни варианти и обосновка на избора}
\label{subsec:choices}

\subsubsection{Място на изчислението}
\label{subsec:where}

Изчислението е справка в таблица, не числено решаване: за дадена комбинация от входни
величини се намира ред и от него се четат до дванадесет числа. Това прави въпроса
„къде да се извърши“ въпрос за разположението на \emph{данните}, а не за изчислителен
ресурс. Разгледани са три варианта. При \textbf{А} клиентът изпраща избора и сървърът
връща готовите числа: една реализация и един източник на истина, но всяка промяна на
кой да е избор изисква обръщение по мрежата, а по К2 това е най-слабият вариант,
защото прави калкулатора зависим от достъпността на услугата. При \textbf{Б} клиентът
получава целия набор веднъж и работи локално, което е бързо, но калкулаторът не работи,
докато услугата не отговори поне веднъж. При \textbf{В} наборът се доставя като
статичен файл заедно с останалите ресурси на страницата.

Приет е вариант В за уеб клиента и вариант Б за мобилния. Наборът е с обем 407{,}7 kB
в един статичен файл, съпоставимо с една снимка в страница, което не оправдава
сложност. Уеб клиентът остава напълно работоспособен без сървър, тоест точно
формулировката на NFR2. Мобилният клиент получава същия набор през \code{GET
/api/loads}, защото при него доставката на нова версия минава през магазин за
приложения и вграждането на данните би означавало, че всяка промяна в таблиците
изисква нова версия на приложението.

\subsubsection{Технология на уеб клиента}
\label{subsec:webtech}

Разгледани са приложение с рамка и стъпка на компилация (React с Vite) и приложение на
стандартен JavaScript без стъпка на компилация. И двата варианта бяха реализирани:
първият съществува в хранилището като \code{web/} и обхваща около 1230 реда, а вторият,
\code{website/}, обхваща 4911 реда JavaScript и CSS плюс 822 реда HTML. Приет е
вторият. Решаващ е К3: той се публикува чрез копиране на директория и се обслужва като
статични файлове от същия сървър, който носи и приложния програмен интерфейс, без нито
една стъпка на подготовка. Второ съображение е дълготрайността: системата е
предназначена да бъде поддържана от конструктивен отдел, а не от екип за уеб
разработка, и всяка зависимост от версии на инструментална верига е бъдещ разход.

\subsubsection{Механизъм на удостоверяване}
\label{subsec:authchoice}

Разгледани са сесия, съхранявана на сървъра (позволява незабавно прекратяване, но
изисква състояние и обръщение към хранилището при всяка заявка), подписан маркер в
\code{localStorage} (еднакво работи навсякъде, но е достъпен от всеки скрипт на
страницата, което противоречи на NFR5) и подписан маркер в бисквитка с признак
\code{httpOnly} (недостъпен от скрипт, но бисквитките не са естественият механизъм за
мобилен клиент). Прието е съчетание: бисквитка за уеб клиента и маркер в заглавие за
мобилния, с обща проверка на сървъра. При регистрация и вход сървърът едновременно
поставя бисквитка \code{wt\_token} с признаци \code{httpOnly} и \code{sameSite=lax} и
връща същия маркер в тялото на отговора. Защитените маршрути приемат и двата източника,
което позволява един приложен програмен интерфейс да обслужва двата клиента, без нито
един от тях да прави компромис с начина, по който съхранява маркера.

\subsubsection{Подход към мобилния клиент}
\label{subsec:mobilechoice}

Разгледани са отделни приложения на естествения за всяка платформа език (пълен достъп
до платформата, но две реализации на един интерфейс, тоест К4 е нула), прогресивно уеб
приложение (няма втора реализация, но не се разпространява през магазин за приложения,
а именно това е очакваният канал) и React Native с Expo. Приет е React Native.
Решаващ е К4: изчислителното ядро \code{lib/loads.js} и построяването на схемата
\code{lib/acsGeometry.js} се използват от уеб прототипа и от мобилното приложение
\emph{без нито един различен знак}, което е проверено чрез побайтово сравнение на
файловете. Така NFR1 е свойство на един файл, а не твърдение, което трябва да се
доказва поотделно за всеки клиент.

\subsubsection{Обобщение}
\label{subsec:choice-summary}

\begin{table}[htbp]
\centering
\caption{Приети архитектурни решения и водещ критерий}
\label{tab:choices}
{\fontsize{10}{12}\selectfont
\begin{tabular}{@{}L{3.5cm}L{6.4cm}L{2.6cm}@{}}
\toprule
\textbf{Въпрос} & \textbf{Прието решение} & \textbf{Водещ критерий} \\
\midrule
Място на изчислението & в клиента; данните вградени за уеб, доставени по мрежата за мобилния & К2 \\
Технология на уеб клиента & стандартен JavaScript без стъпка на компилация & К3 \\
Удостоверяване & подписан маркер, бисквитка \code{httpOnly} за уеб и заглавие \code{Bearer} за мобилния & К1, К4 \\
Мобилен клиент & React Native (Expo) със споделено изчислително ядро & К4 \\
\bottomrule
\end{tabular}}
\end{table}

Разгледаните решения не са равностойни по цена на грешката. Единственото, чието
връщане назад би било скъпо, е мястото на изчислението, защото то определя формата на
данните и договора на интерфейса. Останалите три са локални: смяната на технологията
на уеб клиента засяга една директория, смяната на механизма за удостоверяване засяга
два файла на сървъра, а мобилният клиент може да бъде заменен изцяло, без сървърът да
се промени. Това разпределение е съзнателно и е причина първото решение да бъде
обосновано най-подробно.

\subsection{Обща архитектура}
\label{subsec:arch}

Системата се състои от една услуга и два клиента, които общуват с нея през един и същ
приложен програмен интерфейс \figref{fig:arch}. Услугата изпълнява две несвързани
задачи: обслужва приложния програмен интерфейс и раздава статичните файлове на уеб
клиента. Второто не е техническа необходимост, а решение с последица за
удостоверяването: понеже страницата и интерфейсът се обслужват от един и същ произход,
бисквитката със сесията се изпраща без допълнителна уговорка за споделяне на ресурси
между произходи.

\begin{figure}[htbp]
\centering
\begin{tikzpicture}[node distance=1.0cm and 1.4cm]
  \node[block] (web) {Уеб клиент\\{\scriptsize браузър, статични файлове}};
  \node[block, below=1.1cm of web] (mob) {Мобилен клиент\\{\scriptsize React Native (Expo)}};
  \node[block, right=2.2cm of $(web)!0.5!(mob)$] (api) {Услуга\\{\scriptsize Node.js + Express}};
  \node[accent, right=1.6cm of api] (db) {MongoDB Atlas\\{\scriptsize users, projects, questions}};
  \node[note, above=1.0cm of api] (data) {Товарни таблици\\{\scriptsize статичен набор, 407,7 kB}};

  \draw[line] (web) -- node[above,font=\scriptsize,sloped] {бисквитка} (api);
  \draw[line] (mob) -- node[below,font=\scriptsize,sloped] {Bearer} (api);
  \draw[line] (api) -- (db);
  \draw[line] (data) -- node[right,font=\scriptsize] {/api/loads} (api);
  \draw[line, dashed] (data.west) -- ++(-1.2,0) |- node[above,font=\scriptsize,pos=0.25] {вграден} (web.north);
\end{tikzpicture}
\caption{Обща архитектура: една услуга, два клиента, общ приложен програмен интерфейс}
\label{fig:arch}
\end{figure}

Легенда към \figref{fig:arch}: плътните стрелки са обръщения по мрежата, прекъснатата
стрелка означава, че наборът с товарните таблици се доставя на уеб клиента като част
от статичните файлове, а не по мрежата. Надписите върху стрелките от клиентите към
услугата означават начина, по който всеки от тях носи маркера за удостоверяване.

\subsection{Слоеве и отговорности}
\label{subsec:layers}

Всеки от трите изпълними компонента е разделен на едни и същи три слоя, което прави
съответствията между тях преки:

\begin{itemize}[topsep=2pt,itemsep=2pt,leftmargin=*]
\item \textbf{Слой на данните.} В услугата това е единствената връзка към базата от
      данни заедно с указателите; в клиентите това е достъпът до товарните таблици.
      Този слой не знае нищо за интерфейса.
\item \textbf{Изчислителен слой.} Свеждането на входните величини до ред от таблицата,
      ограничаването на избора до допустимите стойности, прилагането на частните
      коефициенти и формàтирането на числата. Не докосва документа на страницата.
\item \textbf{Слой на представянето.} Съставянето на разметката, свързването на
      събитията и графичната схема. Не знае откъде идват числата.
\end{itemize}

Разделението е проведено докрай в едно отношение и съзнателно не докрай в друго. В
уеб прототипа и в мобилното приложение изчислителният слой е отделен файл
(\code{lib/loads.js}), който е общ за двата. В приетия уеб клиент изчислителният слой
и слоят на представянето живеят в един файл (\code{app.js}), защото извеждането им в
отделен модул би наложило или стъпка на компилация, или система от модули в браузъра,
и двете в противоречие с К3 от §\ref{subsec:criteria}. Разделението там е по
функции, а не по файлове, и е записано като известно ограничение в
§\ref{subsec:webimpl}.

\subsection{Проектиране на базата данни}
\label{subsec:datamodel}

Базата от данни съдържа три колекции \tabref{tab:collections}, свързани в проста
верига: един потребител притежава много проекти, а един проект носи много запитвания
\figref{fig:er}. Товарните таблици \emph{не} се съхраняват в базата: те са едни и същи
за всички потребители, не се променят от системата и се доставят като статичен ресурс
(§\ref{subsec:where}). Това е и причината базата да съдържа само потребителски данни,
а отпадането ѝ да не спира калкулатора.

\begin{figure}[htbp]
\centering
\begin{tikzpicture}[node distance=2.2cm, every node/.append style={font=\scriptsize}]
  \node[block] (u) {\textbf{users}\\ \_id\\ name\\ email\\ passwordHash};
  \node[block, right=3.0cm of u] (p) {\textbf{projects}\\ \_id\\ userId\\ name, tags[]\\ properties[]\\ input, snapshot};
  \node[block, right=3.0cm of p] (q) {\textbf{questions}\\ \_id\\ projectId\\ userId\\ message, status};
  \draw[line] (u) -- node[above] {1 \quad $\ast$} node[below] {userId} (p);
  \draw[line] (p) -- node[above] {1 \quad $\ast$} node[below] {projectId} (q);
\end{tikzpicture}
\caption{Връзки между колекциите: един потребител притежава много проекти, един проект носи много запитвания}
\label{fig:er}
\end{figure}

Легенда към \figref{fig:er}: „1“ означава единичния край на връзката, „$\ast$“ множествения;
надписът под стрелката е полето, което носи връзката. Изтриването на проект изтрива и
неговите запитвания; изтриването се извършва от приложния слой, тъй като документната
база не налага външни ключове.

Изборът на документна база е продиктуван от едно свойство на предметната област:
потребителските свойства на проекта са произволни двойки „ключ: стойност“, чийто набор
не е известен предварително и се различава за всеки клиент. В релационен модел това
изисква или отделна таблица за свойствата, или колона със свободен текст; в документен
модел то е естествено вложено поле.

\begin{table}[htbp]
\centering
\caption{Колекции в базата от данни}
\label{tab:collections}
{\fontsize{10}{12}\selectfont
\begin{tabular}{@{}L{2.4cm}L{6.6cm}L{4.0cm}@{}}
\toprule
\textbf{Колекция} & \textbf{Полета} & \textbf{Указатели} \\
\midrule
\code{users} & \code{name}, \code{email}, \code{passwordHash}, \code{createdAt} & \code{email} (уникален) \\
\code{projects} & \code{userId}, \code{name}, \code{tags[]}, \code{properties[\{key,value\}]}, \code{input\{\}}, \code{snapshot\{\}}, \code{createdAt}, \code{updatedAt} & \code{(userId, updatedAt)}, \code{(userId, tags)} \\
\code{questions} & \code{projectId}, \code{userId}, \code{projectName}, \code{message}, \code{status}, \code{createdAt}, \code{emailStatus} & \code{(projectId, createdAt)} \\
\bottomrule
\end{tabular}}
\end{table}

Проектът пази две различни по предназначение вложени полета. Полето \code{input} е
\emph{състоянието на калкулатора}, тоест деветте величини, от които резултатът може да
бъде възпроизведен изцяло. Полето \code{snapshot} е \emph{обобщение за показване}:
заглавие, мерна единица, водеща стойност и оценка. Дублирането е съзнателно и има
точна причина: списъкът с проекти показва обобщението за десетки проекта
едновременно, а възпроизвеждането му от \code{input} би изисквало клиентът да
изпълни изчислението за всеки ред от списъка. Цената на решението е, че при промяна
в товарните таблици обобщенията остават стари, докато проектът не бъде отворен и
записан отново; това е записано като ограничение в §\ref{subsec:limitations}.

Уникалният указател върху \code{email} не е оптимизация, а ограничение за цялост:
той е единственото място, което не допуска два профила с една и съща поща, и
поражда грешка с код 11000, която слоят на удостоверяването превежда в отговор 409.
Проверка „има ли вече такъв адрес“ преди вмъкването не би заменила указателя, защото
между проверката и вмъкването може да се вклини втора заявка; указателят решава
състезанието, проверката само го прави по-рядко.

Останалите два указателя следват пряко от заявките, които интерфейсът поражда.
Съставният указател \code{(userId, updatedAt)} обслужва списъка с проекти, който
винаги е ограничен до един потребител и по подразбиране е подреден по време на
последна промяна. Съставният указател \code{(userId, tags)} обслужва филтрирането по
етикет. И двата започват с \code{userId}, защото няма нито една заявка, която да чете
проекти на повече от един потребител: обхватът на достъпа е част от условието за
търсене, не се прилага след него.

\subsection{Договор на приложния програмен интерфейс}
\label{subsec:api}

Интерфейсът е разделен на три групи според изискването за удостоверяване
\tabref{tab:api}. Разделението не е само организационно: публичната група е монтирана
\emph{преди} проверката за връзка с базата от данни, а останалите след нея. Това е
техническата реализация на NFR2.

\begin{table}[htbp]
\centering
\caption{Входни точки на приложния програмен интерфейс}
\label{tab:api}
{\fontsize{10}{12}\selectfont
\begin{tabular}{@{}L{1.75cm}L{5.95cm}L{5.3cm}@{}}
\toprule
\textbf{Метод} & \textbf{Път} & \textbf{Предназначение} \\
\midrule
\multicolumn{3}{@{}l}{\textit{Публични, не изискват база от данни}} \\
GET & \code{/api/health} & състояние на услугата и на връзката с базата \\
GET & \code{/api/loads} & пълният набор товарни таблици \\
\addlinespace
\multicolumn{3}{@{}l}{\textit{Удостоверяване (с ограничение на честотата)}} \\
POST & \code{/api/auth/register} & създаване на профил \\
POST & \code{/api/auth/login} & вход \\
POST & \code{/api/auth/logout} & изход \\
GET & \code{/api/auth/me} & текущият профил \\
\addlinespace
\multicolumn{3}{@{}l}{\textit{Изискват удостоверяване, обхват само собствените записи}} \\
GET & \code{/api/projects} & списък с филтри \code{?tag=\&q=\&sort=} \\
GET & \code{/api/projects/meta/tags} & различните етикети на потребителя \\
POST & \code{/api/projects} & създаване \\
GET & \code{/api/projects/:id} & отваряне \\
PUT & \code{/api/projects/:id} & промяна \\
DELETE & \code{/api/projects/:id} & изтриване, заедно със запитванията \\
POST & \code{/api/projects/:id/}\newline\code{support-inquiries} & запитване към инженерния отдел \\
\bottomrule
\end{tabular}}
\end{table}

Договорът е проектиран около две правила, които важат за всяка входна точка.
Първото: обхватът на всяка операция върху проект включва притежателя, тоест условието
за търсене винаги съдържа \code{userId}, а не само \code{\_id}. Това превръща опита за
достъп до чужд проект в отговор 404, а не в 403, и не разкрива дали такъв проект
съществува. Второто: отговорът никога не съдържа вътрешни полета; всеки документ
минава през преобразуване, което изгражда публичното представяне поле по поле, вместо
да премахва полета от вътрешното.

\subsection{Проектиране на изчислителното ядро}
\label{subsec:core}

Изчислителното ядро има едно свойство, което определя цялото му устройство: наборът
от допустими стойности за един избор зависи от другите избори. Височина 9 m допуска
само схема с 2 нива; смяната на типа от катерачна на боулдер стена прави текущата
височина невалидна, защото боулдер стените съществуват само за 4 и 5 m. Ако
интерфейсът позволи да се задържи невалидна комбинация, потребителят вижда съобщение
за липсващ ред там, където би трябвало да види резултат.

Затова ядрото е проектирано около една операция, наречена \term{привеждане}
(\code{clampAndRender}), която се изпълнява след \emph{всяка} промяна на вход и
привежда цялото състояние към най-близкото допустимо \figref{fig:clamp}. Логиката ѝ
е следната: при невалидна стойност се избира стойност по подразбиране за височината,
последната налична схема, най-голямото допустимо разстояние между колоните и така
нататък, и едновременно с това се отбелязва, че тази величина \emph{не е избрана от
потребителя}. Второто е нужно, защото системата показва резултат само когато всички
величини са избрани съзнателно; иначе на екрана би се появила таблица с числа, които
потребителят не е поискал.

\begin{figure}[htbp]
\centering
\begin{tikzpicture}[node distance=0.72cm]
  \node[block] (in) {Промяна на вход};
  \node[block, below=of in] (type) {Тип съвместим ли е\\с височината?};
  \node[danger, right=1.7cm of type] (fixh) {височина по\\подразбиране};
  \node[block, below=of type] (sch) {Схемата налична ли е\\за тази височина?};
  \node[danger, right=1.7cm of sch] (fixs) {последна\\налична схема};
  \node[block, below=of sch] (ax) {$A$ и $X$ в обхвата\\на конфигурацията?};
  \node[danger, right=1.7cm of ax] (fixax) {най-близка\\допустима стойност};
  \node[block, below=of ax] (all) {Всички величини\\избрани съзнателно?};
  \node[note, right=1.7cm of all] (empty) {показва се покана\\за избор};
  \node[accent, below=of all] (out) {Извеждане на\\резултата};

  \draw[line] (in) -- (type);
  \draw[line] (type) -- node[right,font=\scriptsize,pos=0.25] {да} (sch);
  \draw[line] (type) -- node[above,font=\scriptsize] {не} (fixh);
  \draw[line] (fixh.south) |- (sch.east);
  \draw[line] (sch) -- node[right,font=\scriptsize,pos=0.25] {да} (ax);
  \draw[line] (sch) -- node[above,font=\scriptsize] {не} (fixs);
  \draw[line] (fixs.south) |- (ax.east);
  \draw[line] (ax) -- node[right,font=\scriptsize,pos=0.25] {да} (all);
  \draw[line] (ax) -- node[above,font=\scriptsize] {не} (fixax);
  \draw[line] (fixax.south) |- (all.east);
  \draw[line] (all) -- node[right,font=\scriptsize,pos=0.25] {да} (out);
  \draw[line] (all) -- node[above,font=\scriptsize] {не} (empty);
\end{tikzpicture}
\caption{Привеждане на състоянието към допустима комбинация след промяна на вход}
\label{fig:clamp}
\end{figure}

Легенда към \figref{fig:clamp}: правоъгълниците със заоблени ъгли са проверки,
светлочервените блокове са поправки на състоянието, а зеленият блок е изходът.
Всяка поправка отбелязва съответната величина като неизбрана, което е причината
потребителят да вижда поканата за избор вместо резултат, докато не потвърди
поправената стойност.

Второ проектно решение в ядрото е отделянето на \term{избор} от \term{преобразуване}.
Смяната на мерната система не преобразува числа, а превключва към другия набор в
същия ред от таблицата, точно както е обосновано в §\ref{subsec:tables}. Единственото
преобразуване, което системата извършва, е върху \emph{дължините} за показване
(метри към футове), защото те са входни величини на потребителя, а не стойности от
еталона.

\subsection{Проектиране на потребителския интерфейс}
\label{subsec:uidesign}

Интерфейсът на калкулатора е изграден като два дяла: неподвижен дял с входните
величини и дял с резултата. На настолен екран те стоят един до друг, като дялът с
входните величини е прилепен и има собствено превъртане, за да остане видим, докато
потребителят разглежда резултата. Под определена ширина двата дяла се подреждат един
под друг.

Разположението е зададено с решетка от две колони, при която първата колона има
променлива, но ограничена ширина, а втората заема остатъка. Точките на пречупване са
изведени не от размерите на конкретни устройства, а от момента, в който съдържанието
престава да се побира \tabref{tab:breakpoints}.

\begin{table}[htbp]
\centering
\caption{Точки на пречупване и промяната, която въвеждат}
\label{tab:breakpoints}
{\fontsize{10}{12}\selectfont
\begin{tabular}{@{}L{2.4cm}L{10.6cm}@{}}
\toprule
\textbf{Ширина} & \textbf{Промяна} \\
\midrule
$\le$ 520 px & намалени външни отстъпи; списъкът с проекти е една колона \\
$\le$ 600 px & сгъстена таблица с товарите; избор на ниво в две колони \\
$\le$ 640 px & лентата с филтри става решетка етикет/поле \\
$\le$ 720 px & подравняване на абзаците вляво; лентата с инструменти престава да е прилепена; навигацията в документацията става две колони \\
$\le$ 760 px & двете схеми на закрепване една под друга; графиката с ограничена височина \\
$\le$ 960 px & входните величини и резултатът един под друг; заглавната лента става решетка от два реда \\
$\le$ 1300 px & намален набор от размери на шрифта \\
$\le$ 1450 px & съкратена основна навигация \\
\bottomrule
\end{tabular}}
\end{table}

Ключово проектно правило, извлечено от изискването NFR4, е следното: \emph{никой
елемент не разширява страницата}. Съдържание, което по естеството си е по-широко от
екрана (таблица с шест колони, чертеж в мащаб 1:1, обзорен лист с детайли за
закрепване), се поставя в собствен превъртан съд, а не се оставя да разшири
документа. Разликата между двете е съществена за използваемостта: в първия случай
потребителят превърта таблицата и вижда цялата, във втория цялата страница се измества
настрани, а при налична забрана за хоризонтално превъртане на документа съдържанието
просто остава недостъпно.

\subsection{Проектиране на графичната схема}
\label{subsec:schematic}

Схемата на съоръжението не е изображение, а чертеж, изчислен от текущия избор:
височина, разстояние между колоните, наклон и височини на нивата на закрепване.
Проектирана е като векторна графика, съставяна в клиента, с три следствия. Първо,
схемата е точна при всяко увеличение, което има значение, защото тя носи размерни
линии. Второ, отделните ѝ елементи са част от документа и могат да носят събития,
което позволява точките на закрепване да бъдат посочени и да отворят съответния
типов детайл. Трето, схемата се променя със същата операция, която преизчислява
таблицата, така че двете не могат да се разминат.

Изборът на съотношение и на видимо поле е направен така, че чертежът да запълва
наличната ширина, а височината му да следва съдържанието. При тесен екран схемата
получава ограничена височина и собствено превъртане с преместване и увеличение, вместо
да бъде смалена до нечетливост.
